\section{Related Work}

\textbf{Autonomous AI Research Agents.}
Recent AI scientist agents aim to automate large parts of the research pipeline. \textit{The AI Scientist}~\citep{lu2024aiscientist} and \textit{The AI Scientist-v2}~\citep{yamada2025aiscientistv2} are early systems that attempt this process by autonomously generating ideas, writing code, running experiments, and drafting papers. \textit{Agent Laboratory}~\citep{schmidgall2025agentlab} structures this process into literature review, experimentation, and report writing, with optional human feedback. Beyond papers, \textit{autoresearch}~\citep{karpathy2026autoresearch} provides a simple open-source loop where an agent edits code, runs short real experiments, and keeps changes only when metrics improve. These systems show that LLM agents can execute most parts of research workflows. However, they do not directly test whether LLMs can judge whether a research proposal (including the core idea, related work context, and experimental design) is worth pursuing \emph{before} execution begins. Like human researchers triaging projects, this first-gate decision determines whether time and compute should be invested.

\textbf{Evaluating Research Agents and Idea Quality.}
% Recent studies suggest that LLM-generated ideas can be novel while being weaker on feasibility~\citep{si2024llmideas} and this gap can widen after implementation~\citep{si2025ideationexecution}. At the system level, most current benchmarks focus on evaluating execution~\citep{chan2024mle,starace2025paperbench,wu2025innovatorbench,lupidi2026airs}. Additionally, a few idea-evaluation benchmarks target dimensions such as impact and novelty~\citep{jiang2026hindsight, rinobench2026}, or model scientific taste from community-level signals (publication venue and citation outcomes)~\citep{tong2026scientifictaste}. These directions are important, but their targets differ from ours. In contrast, SoundnessBench targets a narrower upstream question: whether a research proposal is methodologically sound before results exist and whether it is a good test of the claimed idea.
\fh{Recent studies suggest that LLM-generated ideas can be novel while being weaker on feasibility~\citep{si2024llmideas}, and that the gap between proposal-stage promise and post-execution outcomes can widen after implementation~\citep{si2025ideationexecution}. We take this implication seriously: proposal-only judgment may be intrinsically limited because some weaknesses or strengths become visible only after experiments are run. SoundnessBench is therefore not intended to predict final research impact or full-paper acceptance. Instead, it asks a narrower diagnostic question: how much reviewer-assigned methodological soundness is recoverable from proposal-stage evidence, and whether LLMs can act as stable first-gate critics under that limitation. At the system level, most current benchmarks focus on evaluating execution~\citep{chan2024mle,starace2025paperbench,wu2025innovatorbench,lupidi2026airs}. Additionally, a few idea-evaluation benchmarks target dimensions such as impact and novelty~\citep{jiang2026hindsight, rinobench2026}, or model scientific taste from community-level signals such as publication venue and citation outcomes~\citep{tong2026scientifictaste}. These directions are important, but their targets differ from ours.}

\textbf{LLMs as Peer Reviewers.}
Existing work studies LLMs as reviewers of completed papers. LLM assistance can improve review clarity and usefulness in human review workflows~\citep{thakkar2025llmfeedback}. Recent benchmarks also study the behavior and quality of LLM-based reviewing systems at scale~\citep{omnireview2026, isyourpaperreviewed2025}. They study how models produce critiques and how closely their judgments track human assessments. However, this literature evaluates \emph{post hoc} review, where the paper, results, and framing are already available. SoundnessBench instead studies \emph{pre-execution} judgment from the research proposal alone.

\textbf{LLM Sycophancy and Fragile Judgment.}
Our work is also connected to studies of LLM sycophancy and fragile judgment. Previous work demonstrates that LLMs often align with user framing rather than ground truth, which can make them validate plausible but weak claims~\citep{sharma2024sycophancy}. Other work shows that attempts to suppress this behavior can produce unstable skepticism or excessive refusal~\citep{causalt5k}. Prompt-framing studies further show that small syntactic changes can flip model decisions even when content is unchanged~\citep{syntacticfragility}. These findings suggest that scientific criticism is not only a knowledge problem, but also a robustness problem. SoundnessBench tests both.